\documentclass[11pt,a4paper]{article}
\usepackage[margin=1in]{geometry}
\usepackage{hyperref}
\usepackage{enumitem}
\usepackage{tcolorbox}
\usepackage{booktabs}

\usepackage{xcolor}
\definecolor{listbackgroundcolorlight}{rgb}{0.91,0.92,0.94}
\definecolor{keywordColor}{HTML}{800080}
\definecolor{commentColor}{HTML}{008000}
\definecolor{colorEntityBack}{rgb}{0.01,0.01,0.4}
\definecolor{colorAbstractionLayer}{rgb}{0.01,0.3,0.7}
\definecolor{colorPolicyBack}{rgb}{0.91,0.94,0.94}
\definecolor{labelcolor}{cmyk}{0.22,0.10,0.10,0.10}
\definecolor{colorProtocol}{rgb}{0.6,0.2,0.4}
%% \usepackage{crayola}
\definecolor{colorOrg}{gray}{0.85}

\usepackage[formats]{listings}
\lstdefineformat{BSPLformat}{->=$\textcolor{keywordColor}{\mapsto}$}
\lstset{format=BSPLformat}
\lstdefinelanguage{BSPL}
{
  morekeywords={roles, parameters, private, out, in, nil, key},
  otherkeywords={->,:,],[},
  morecomment=[l]{//}
}

% Code listing style
\lstset{
    basicstyle=\ttfamily\small,
    keywordstyle=\color{blue},
    stringstyle=\color{red},
    commentstyle=\color{gray},
    breaklines=true,
    frame=single,
    numbers=left,
    numberstyle=\tiny\color{gray}
}
\lstset{showspaces=false}
\lstset{keepspaces=true}

\title{\textbf{BSPL FastAPI Server \& Wrapper} \\ 
\large Quickstart Guide \& API Reference}
\author{LLM BSPL Adapter Project}
\date{\today}

\begin{document}

\maketitle

\begin{abstract}
This document provides a comprehensive guide to the BSPL FastAPI server and wrapper functions. It serves as both a quickstart tutorial and a complete API reference, covering server setup, configuration, core concepts, and detailed endpoint documentation with practical examples.
\end{abstract}

\tableofcontents
\newpage

%========================================
\section{Introduction}
%========================================

The BSPL FastAPI Server provides a REST API interface to the Business Protocol Specification Language (BSPL) adapter, enabling protocol-driven multi-agent systems to interact via HTTP endpoints. The system consists of:

\begin{itemize}
    \item \textbf{FastAPI Server} (\texttt{server.py}): HTTP REST API exposing protocol operations
    \item \textbf{BSPLConcreteWrapper} (\texttt{bspl\_concrete\_wrapper.py}): Bridge layer between API and core adapter
    \item \textbf{Core Adapter} (\texttt{core.py}): BSPL protocol execution engine
\end{itemize}

\subsection{Key Features}

\begin{itemize}[leftmargin=*]
    \item \textbf{Stateful Protocol Management}: Maintain protocol state across API calls
    \item \textbf{Message Simulation}: Preview message effects without state mutation (dry-run mode)
    \item \textbf{Idempotency Support}: Safe message retries with idempotency keys
    \item \textbf{Policy-Based Authorization}: Pluggable validators for role-based access control
    \item \textbf{LLM-Ready}: OpenAI-style function definitions for AI agent integration
    \item \textbf{Cross-Platform}: Windows (PowerShell) and Linux (Bash) support
\end{itemize}

%========================================
\section{Quickstart Guide}
%========================================

\subsection{Prerequisites}

\begin{lstlisting}[language=bash]
# Install dependencies
pip install fastapi uvicorn pydantic

# Ensure BSPL package is installed
pip install -e .
\end{lstlisting}

\subsection{Starting the Server}

\begin{lstlisting}[language=bash]
# Start server on default port 8000
python -m uvicorn bspl.adapter.server:app --host 0.0.0.0 --port 8000

# With auto-reload for development
python -m uvicorn bspl.adapter.server:app --reload
\end{lstlisting}

\subsection{Verifying Server Status}

\begin{lstlisting}[language=bash]
# Check server health
curl http://localhost:8000/health

# View server info and available endpoints
curl http://localhost:8000/
\end{lstlisting}

%========================================
\section{Exposed Functions}
%========================================

\subsection{Configure the Adapter}

Before using protocol operations, configure the adapter with your protocol definition and agent mappings.

\begin{tcolorbox}[colback=blue!5,colframe=blue!40!black,title=Configuration Requirements]
\textbf{Important:} For now, protocol objects cannot be serialized to JSON. Configuration must be done programmatically via the Python API, not through the HTTP endpoint. I'm putting this on the long to-do if I get time as that would make the LLM-adapter interface cleaner.
\end{tcolorbox}

\begin{lstlisting}[language=Python]
# In your application code
from bspl.adapter.server import configure_adapter
import bspl

# Load protocol
purchase_protocol = bspl.load_file(
    "scenarios/purchase/purchase.bspl"
).protocols["Purchase"]

# Define agent network topology
agents = {
    "Buyer": [("127.0.0.1", 9000)],
    "Seller": [("127.0.0.1", 9001)],
    "Shipper": [("127.0.0.1", 9002)],
    "my_adapter": [("127.0.0.1", 8000)]
}

# Map protocol roles to agents
systems = {
    "purchase_system": {
        "protocol": purchase_protocol,
        "roles": {
            purchase_protocol.roles["Buyer"]: "my_adapter",
            purchase_protocol.roles["Seller"]: "Seller",
            purchase_protocol.roles["Shipper"]: "Shipper"
        }
    }
}

# Configure the adapter
configure_adapter(
    systems=systems,
    agents=agents,
    name="my_adapter"
)
\end{lstlisting}

\subsection{Query Available Functions}

For LLM integration, retrieve function definitions:

\begin{lstlisting}[language=bash]
curl http://localhost:8000/llm_functions
\end{lstlisting}

\textbf{Response:} Returns OpenAI-style function schemas for all available operations.

\subsection{Check Role State}

Query the current state for a specific role:

\begin{lstlisting}[language=bash]
curl -X POST http://localhost:8000/get_role_state \
  -H "Content-Type: application/json" \
  -d '{"role_id": "Buyer"}'
\end{lstlisting}

\textbf{Response includes:}
\begin{itemize}
    \item Message history per protocol system
    \item Variable bindings in each context
    \item Enabled (pending) messages
    \item Active protocols and roles
\end{itemize}

\subsection{Step 4: List Enabled Messages}

See what messages can currently be sent:

\begin{lstlisting}[language=python]
curl -X POST http://localhost:8000/list_enabled_messages \
  -H "Content-Type: application/json" \
  -d '{"role_id": "Buyer"}'
\end{lstlisting}

\subsection{Step 5: Simulate a Message (Dry Run)}

This can be used for few-shot training. Basically can preview message effects without committing:

\begin{lstlisting}[language=bash,showspaces=false]
curl -X POST http://localhost:8000/simulate_message \
  -H "Content-Type: application/json" \
  -d '{
    "schema": "Purchase.rfq",
    "payload": {
      "ID": "order-123",
      "item": "laptop",
      "quantity": 5
    },
    "system": "purchase_system"
  }'
\end{lstlisting}

\textbf{Response shows:}
\begin{itemize}
    \item Messages that would be enabled
    \item Messages that would be disabled
    \item Potential state observations
\end{itemize}

\subsection{Step 6: Send a Real Message}

Execute the protocol transition:

\begin{lstlisting}[language=bash]
curl -X POST http://localhost:8000/send_message \
  -H "Content-Type: application/json" \
  -d '{
    "schema": "Purchase.rfq",
    "payload": {
      "ID": "order-123",
      "item": "laptop",
      "quantity": 5
    },
    "system": "purchase_system",
    "dry_run": false,
    "idempotency_key": "rfq-order-123-attempt-1"
  }'
\end{lstlisting}

\textbf{Idempotency:} Using the same \texttt{idempotency\_key} will return the cached result instead of re-sending.

\subsection{Step 7: Receive External Messages}

Ingest messages from other agents:

\begin{lstlisting}[language=bash]
curl -X POST http://localhost:8000/receive \
  -H "Content-Type: application/json" \
  -d '{
    "schema": "Purchase.quote",
    "payload": {
      "ID": "order-123",
      "price": 5000,
      "delivery_date": "2025-11-01"
    },
    "from": "Seller"
  }'
\end{lstlisting}

\subsection{Step 8: Update Local State (Advanced)}

Manually inject state updates (dangerous!):

\begin{lstlisting}[language=bash]
curl -X POST http://localhost:8000/update_local_state \
  -H "Content-Type: application/json" \
  -d '{
    "role_id": "Buyer",
    "updates": {
      "inventory_level": 10,
      "last_order_id": "order-123"
    }
  }'
\end{lstlisting}

\subsection{Step 9: Abort Protocol (Emergency)}

Stop protocol execution for a role:

\begin{lstlisting}[language=bash]
curl -X POST http://localhost:8000/abort_protocol \
  -H "Content-Type: application/json" \
  -d '{
    "role_id": "Buyer",
    "reason": "Emergency shutdown requested"
  }'
\end{lstlisting}

%========================================
\section{API Reference}
%========================================

\subsection{Informational Endpoints}

\subsubsection{GET /}
\textbf{Description:} Server status and endpoint listing

\textbf{Response:}
\begin{lstlisting}[language=python]
{
  "service": "BSPL Adapter Wrapper API",
  "version": "1.0",
  "status": "running",
  "adapter_configured": true,
  "endpoints": { ... }
}
\end{lstlisting}

\subsubsection{GET /health}
\textbf{Description:} Health check endpoint

\textbf{Response:}
\begin{lstlisting}[language=python]
{
  "status": "healthy",
  "adapter_configured": true
}
\end{lstlisting}

\subsubsection{GET /llm\_functions}
\textbf{Description:} OpenAI-style function definitions for LLM integration

\textbf{Response:}
\begin{lstlisting}[language=python]
{
  "functions": [
    {
      "name": "send_message",
      "description": "Send a BSPL message...",
      "parameters": { ... }
    },
    ...
  ]
}
\end{lstlisting}

\subsection{Configuration Endpoints}

\subsubsection{POST /configure}
\textbf{Description:} Configure adapter with systems and agents

\textbf{Note:} Limited to simple configurations without protocol objects. For full setup, use programmatic configuration.

\textbf{Request Body:}
\begin{lstlisting}[language=python]
{
  "systems": { ... },
  "agents": { ... },
  "name": "adapter_name"
}
\end{lstlisting}

\subsection{State Query Endpoints}

\subsubsection{POST /get\_role\_state}
\textbf{Description:} Retrieve complete state for a role

\textbf{Request:}
\begin{lstlisting}[language=python]
{
  "role_id": "Buyer"
}
\end{lstlisting}

\textbf{Response:}
\begin{lstlisting}[language=python]
{
  "role_id": "Buyer",
  "state": {
    "history": {
      "system_id": {
        "bindings": { ... },
        "messages": [ ... ],
        "subcontexts": [ ... ]
      }
    },
    "enabled_messages": [ ... ],
    "roles": ["Buyer", "Seller"],
    "protocols": ["Purchase"]
  }
}
\end{lstlisting}

\subsubsection{POST /list\_pending\_messages}
\textbf{Description:} List currently enabled messages

\textbf{Request:}
\begin{lstlisting}[language=python]
{
  "role_id": "Buyer"
}
\end{lstlisting}

\textbf{Response:}
\begin{lstlisting}[language=python]
{
  "pending": [
    {
      "schema": "Purchase.rfq",
      "payload": { ... },
      "sender": "Buyer",
      "system": "purchase_system",
      "key": "..."
    }
  ]
}
\end{lstlisting}

\subsection{Message Operation Endpoints}

\subsubsection{POST /simulate\_message}
\textbf{Description:} Preview message effects without state mutation

\textbf{Request:}
\begin{lstlisting}[language=python]
{
  "schema": "Purchase.rfq",
  "payload": {
    "ID": "order-123",
    "item": "laptop",
    "quantity": 5
  },
  "system": "purchase_system"
}
\end{lstlisting}

\textbf{Response:}
\begin{lstlisting}[language=python]
{
  "predicted": {
    "added": [ ... ],
    "removed": [ ... ],
    "observations": { ... }
  },
  "warnings": null
}
\end{lstlisting}

\subsubsection{POST /send\_message}
\textbf{Description:} Send a protocol message

\textbf{Request:}
\begin{lstlisting}[language=python]
{
  "schema": "Purchase.rfq",
  "payload": {
    "ID": "order-123",
    "item": "laptop",
    "quantity": 5
  },
  "system": "purchase_system",
  "dry_run": false,
  "idempotency_key": "optional-unique-key"
}
\end{lstlisting}

\textbf{Response:}
\begin{lstlisting}[language=python]
{
  "message_id": "uuid-here",
  "status": "sent",
  "details": {
    "schema": "Purchase.rfq"
  },
  "simulated": false
}
\end{lstlisting}

\textbf{Parameters:}
\begin{itemize}
    \item \texttt{schema} (string, required): Qualified message schema name
    \item \texttt{payload} (object, required): Message content
    \item \texttt{system} (string, optional): System identifier
    \item \texttt{dry\_run} (boolean, default: false): Simulate only
    \item \texttt{idempotency\_key} (string, optional): Prevent duplicate sends
\end{itemize}

\subsubsection{POST /receive}
\textbf{Description:} Ingest external message into adapter

\textbf{Request:} Accepts arbitrary JSON representing received message

\textbf{Response:}
\begin{lstlisting}[language=python]
{
  "status": "ok"
}
\end{lstlisting}

\subsection{Advanced Control Endpoints}

\subsubsection{POST /update\_local\_state}
\textbf{Description:} Manually update role state (powerful operation)

\textbf{Warning:} Should be policy-protected in production

\textbf{Request:}
\begin{lstlisting}[language=python]
{
  "role_id": "Buyer",
  "updates": {
    "custom_key": "custom_value"
  }
}
\end{lstlisting}

\subsubsection{POST /abort\_protocol}
\textbf{Description:} Stop protocol execution for a role

\textbf{Request:}
\begin{lstlisting}[language=python]
{
  "role_id": "Buyer",
  "reason": "Emergency shutdown"
}
\end{lstlisting}

\textbf{Response:}
\begin{lstlisting}[language=python]
{
  "aborted": true,
  "reason": "Emergency shutdown"
}
\end{lstlisting}

%========================================
\section{Wrapper API Reference}
%========================================

\subsection{BSPLConcreteWrapper Class}

The wrapper provides a synchronous interface to the asynchronous adapter.

\subsubsection{Constructor}

\begin{lstlisting}[language=Python]
BSPLConcreteWrapper(
    adapter_instance: Adapter,
    policy_validator: Optional[Callable] = None,
    idempotency_store: Optional[Dict] = None
)
\end{lstlisting}

\textbf{Parameters:}
\begin{itemize}
    \item \texttt{adapter\_instance}: Core BSPL adapter instance
    \item \texttt{policy\_validator}: Function \texttt{(role\_id, action, params) -> bool}
    \item \texttt{idempotency\_store}: Optional external cache for idempotency
\end{itemize}

\subsection{Public Methods}

\subsubsection{simulate\_message()}

\begin{lstlisting}[language=Python]
def simulate_message(
    self,
    payload: SendMessageInput
) -> SimulationResult
\end{lstlisting}

\textbf{Description:} Non-destructive message effect prediction

\textbf{Returns:} \texttt{SimulationResult} with predicted state changes

\subsubsection{send\_message()}

\begin{lstlisting}[language=Python]
def send_message(
    self,
    payload: SendMessageInput
) -> SendMessageResult
\end{lstlisting}

\textbf{Description:} Execute protocol message send

\textbf{Features:}
\begin{itemize}
    \item Idempotency key checking
    \item Dry-run mode support
    \item Policy validation
\end{itemize}

\subsubsection{receive()}

\begin{lstlisting}[language=Python]
def receive(
    self,
    data: Dict[str, Any]
) -> None
\end{lstlisting}

\textbf{Description:} Ingest external message into adapter

\subsubsection{get\_role\_state()}

\begin{lstlisting}[language=Python]
def get_role_state(
    self,
    role_id: str
) -> RoleState
\end{lstlisting}

\textbf{Description:} Extract complete state view for role

\textbf{Returns:}
\begin{itemize}
    \item Message history per system
    \item Variable bindings
    \item Enabled messages
    \item Active protocols and roles
\end{itemize}

\subsubsection{list\_pending\_messages()}

\begin{lstlisting}[language=Python]
def list_pending_messages(
    self,
    role_id: str
) -> ListPendingMessagesResult
\end{lstlisting}

\textbf{Description:} Get currently enabled messages

\subsubsection{update\_local\_state()}

\begin{lstlisting}[language=Python]
def update_local_state(
    self,
    payload: UpdateLocalStateInput
) -> Dict[str, Any]
\end{lstlisting}

\textbf{Description:} Apply manual state updates (policy-protected)

\subsubsection{abort\_protocol()}

\begin{lstlisting}[language=Python]
def abort_protocol(
    self,
    role_id: str,
    reason: Optional[str] = None
) -> AbortResult
\end{lstlisting}

\textbf{Description:} Emergency protocol termination

%========================================
\section{Error Handling}
%========================================

\subsection{Exception Hierarchy}

\begin{itemize}
    \item \texttt{AdapterError} (base class)
    \begin{itemize}
        \item \texttt{TransientError} (HTTP 502): Temporary failures, retry possible
        \item \texttt{SemanticError} (HTTP 400): Protocol/semantic violations
        \item \texttt{PolicyError} (HTTP 403): Authorization failures
    \end{itemize}
\end{itemize}

\subsection{HTTP Status Codes}

\begin{center}
\begin{tabular}{ll}
\toprule
\textbf{Code} & \textbf{Meaning} \\
\midrule
200 & Success \\
400 & Bad request (semantic error) \\
403 & Forbidden (policy violation) \\
422 & Validation error (malformed request) \\
500 & Internal server error \\
502 & Transient error (retry possible) \\
\bottomrule
\end{tabular}
\end{center}

%========================================
\section{Integration Examples}
%========================================

\subsection{Python Client Example}

\begin{lstlisting}[language=Python]
import requests

BASE_URL = "http://localhost:8000"

# Check health
response = requests.get(f"{BASE_URL}/health")
print(response.json())

# Send a message
payload = {
    "schema": "Purchase.rfq",
    "payload": {
        "ID": "order-456",
        "item": "smartphone",
        "quantity": 10
    },
    "system": "purchase_system",
    "idempotency_key": "rfq-order-456"
}
response = requests.post(
    f"{BASE_URL}/send_message",
    json=payload
)
print(response.json())
\end{lstlisting}

\subsection{LLM Integration (Claude/Anthropic)}

\begin{lstlisting}[language=Python]
import anthropic
import requests
import json

# Get function definitions
functions = requests.get(
    "http://localhost:8000/llm_functions"
).json()["functions"]

# Convert to Anthropic tool format
tools = [
    {
        "name": f["name"],
        "description": f["description"],
        "input_schema": f["parameters"]
    }
    for f in functions
]

# Use with Claude API
client = anthropic.Anthropic(api_key="your-api-key")

message = client.messages.create(
    model="claude-3-5-sonnet-20241022",
    max_tokens=1024,
    tools=tools,
    messages=[{
        "role": "user",
        "content": "Send an RFQ for 5 laptops"
    }]
)

# Execute tool calls
for content_block in message.content:
    if content_block.type == "tool_use":
        tool_name = content_block.name
        tool_input = content_block.input
        
        # Call FastAPI endpoint
        result = requests.post(
            f"http://localhost:8000/{tool_name}",
            json=tool_input
        )
        
        print(f"Tool {tool_name} result: {result.json()}")
\end{lstlisting}

%========================================
\section{Troubleshooting}
%========================================

\subsection{Common Issues}

\subsubsection{Server Not Starting}

\textbf{Symptom:} \texttt{ModuleNotFoundError} or import errors

\textbf{Solution:}
\begin{lstlisting}[language=bash]
# Ensure BSPL package is installed
pip install -e .

# Check Python path
export PYTHONPATH="${PYTHONPATH}:/path/to/bspl/src"
\end{lstlisting}

\subsubsection{Adapter Not Configured Error}

\textbf{Symptom:} HTTP 500 with "Adapter not configured" message

\textbf{Solution:} Must configure adapter programmatically before using protocol endpoints. See Section 3.1.

\subsubsection{Circular Reference in State}

\textbf{Symptom:} State query returns simplified response with error message

\textbf{Solution:} This is expected for complex protocol states. Use specific endpoints like \texttt{list\_pending\_messages} for targeted queries.

\subsubsection{Message Schema Not Found}

\textbf{Symptom:} \texttt{SemanticError: Schema 'X' not found}

\textbf{Solution:} Check protocol is loaded and schema name uses qualified format: \texttt{ProtocolName.messageName}

\end{document}
